
\chapter{2017年北京卷（文）}

\section{单选题}

\begin{problem}
已知全集 \(U=\R\)，集合 \(A=\{x\mid x<-2\text{ 或 }x>2\}\)，则 \(C_UA=\)（\quad）.

\choices
    {\(\left(-2,2\right)\)}
    {\(\left(-\infty,-2\right)\cup\left(2,+\infty\right)\)}
    {\(\left[-2,2\right]\)}
    {\(\left(-\infty,-2\right]\cup\left[2,+\infty\right)\)}

\end{problem}

\begin{answer}
C
\end{answer}

\begin{solution}
集合 \(A\) 在实数集 \(U=\R\) 中的补集为 \(C_UA=\left[-2,2\right]\)，故选 C.
\end{solution}

\begin{problem}
若复数 \((1-\i)(a+\i)\) 在复平面内对应的点在第二象限，则实数 \(a\) 的取值范围是（\quad）.

\choices
    {\(\left(-\infty,1\right)\)}
    {\(\left(-\infty,-1\right)\)}
    {\(\left(1,+\infty\right)\)}
    {\(\left(-1,+\infty\right)\)}

\end{problem}

\begin{answer}
B
\end{answer}

\begin{solution}
由题意，\((1-\i)(a+\i)=(a+1)+(1-a)\i\). 对应点在第二象限，所以实部小于 \(0\)，虚部大于 \(0\)，即 \(a+1<0\)，\(1-a>0\)，解得 \(a<-1\). 故选 B.
\end{solution}

\begin{problem}
执行如图所示的程序框图，输出的 \(s\) 值为

\begin{center}
\bitmapinclude[max width=6.0cm,max totalheight=6.5cm]{img/2017/beijing_liberal/q3_fig1.png}
\end{center}

\choices
    {\(2\)}
    {\(\dfrac32\)}
    {\(\dfrac53\)}
    {\(\dfrac85\)}

\end{problem}

\begin{answer}
C
\end{answer}

\begin{solution}
按程序框图依次计算：\(k=0\)，\(s=1\). 当 \(k<3\) 时，执行 \(k=1\)，\(s=\frac{s+1}{s}=2\). 再次循环，\(k=2\)，\(s=\frac{2+1}{2}=\frac32\). 第三次循环，\(k=3\)，\(s=\frac{\frac32+1}{\frac32}=\frac53\). 此时 \(k=3\)，不满足 \(k<3\)，输出 \(s=\dfrac53\). 故选 C.
\end{solution}

\begin{problem}
若 \(x\)，\(y\) 满足 \(x\le3\)，\(x+y\ge2\)，\(y\le x\)，则 \(x+2y\) 的最大值为（\quad）.

\choices
    {\(1\)}
    {\(3\)}
    {\(5\)}
    {\(9\)}

\end{problem}

\begin{answer}
D
\end{answer}

\begin{solution}
由约束条件，目标函数 \(x+2y\) 在可行域的顶点处取到最大值. 可行域的右上顶点为 \((3,3)\)，代入得 \(x+2y=3+2\times3=9\)，故选 D.
\end{solution}

\begin{problem}
已知函数 \(f(x)=3^x-\left(\dfrac13\right)^x\)，则 \(f(x)\)（\quad）.

\choices
    {是偶函数，且在 \(\R\) 上是增函数}
    {是奇函数，且在 \(\R\) 上是增函数}
    {是偶函数，且在 \(\R\) 上是减函数}
    {是奇函数，且在 \(\R\) 上是增函数}

\end{problem}

\begin{answer}
B
\end{answer}

\begin{solution}
因为 \(f(-x)=3^{-x}-\left(\dfrac13\right)^{-x}=-\left(3^x-\left(\dfrac13\right)^x\right)=-f(x)\)，所以 \(f(x)\) 是奇函数. 又因为 \(3^x\) 在 \(\R\) 上单调递增，而 \(\left(\dfrac13\right)^x=3^{-x}\) 在 \(\R\) 上单调递减，所以 \(f(x)=3^x-3^{-x}\) 在 \(\R\) 上单调递增. 故选 B.
\end{solution}

\begin{problem}
某三棱锥的三视图如图所示，则该三棱锥的体积为

\begin{center}
\bitmapinclude[max width=6.0cm,max totalheight=4.8cm]{img/2017/beijing_liberal/q6_fig1.png}
\end{center}

\choices
    {\(60\)}
    {\(30\)}
    {\(20\)}
    {\(10\)}

\end{problem}

\begin{answer}
D
\end{answer}

\begin{solution}
由三视图可得，该三棱锥的底面是直角三角形，底面两直角边长分别为 \(5\) 和 \(3\)，高为 \(4\). 所以体积为 \(V=\frac13\times\frac12\times5\times3\times4=10\)，故选 D.
\end{solution}

\begin{problem}
设 \(m\)，\(n\) 为非零向量，则“存在负数 \(\lambda\)，使得 \(m=\lambda n\)”是“\(m\cdot n<0\)”的（\quad）.

\choices
    {充分而不必要条件}
    {必要而不充分条件}
    {充分必要条件}
    {既不充分也不必要条件}

\end{problem}

\begin{answer}
A
\end{answer}

\begin{solution}
若存在负数 \(\lambda\)，使得 \(m=\lambda n\)，则 \(m\cdot n=\lambda |n|^2<0\)，所以“存在负数 \(\lambda\)，使得 \(m=\lambda n\)”是“\(m\cdot n<0\)”的充分条件. 但当 \(m\cdot n<0\) 时，\(m\) 与 \(n\) 未必共线，因此不是必要条件. 故选 A.
\end{solution}

\begin{problem}
根据有关资料，围棋状态空间复杂度的上限 \(M\) 约为 \(3^{361}\)，而可观测宇宙中普通物质的原子总数 \(N\) 约为 \(10^{80}\). 则下列各数中与 \(\dfrac{M}{N}\) 最接近的是（\quad）.

（参考数据：\(\lg 3\approx0.48\)）

\choices
    {\(10^{33}\)}
    {\(10^{53}\)}
    {\(10^{73}\)}
    {\(10^{93}\)}

\end{problem}

\begin{answer}
D
\end{answer}

\begin{solution}
因为 \(\log_{10}\frac{M}{N}=361\log_{10}3-80\approx361\times0.48-80=93.28\)，所以 \(\frac{M}{N}\approx10^{93}\). 故选 D.
\end{solution}
\section{填空题}

\begin{problem}
在平面直角坐标系 \(xOy\) 中，角 \(\alpha\) 与角 \(\beta\) 均以 \(Ox\) 为始边，它们的终边关于 \(y\) 轴对称. 若 \(\sin\alpha=\dfrac13\)，则 \(\sin\beta=\fillinblank{}\).

\end{problem}

\begin{answer}
\(\dfrac13\)
\end{answer}

\begin{solution}
角 \(\alpha\) 与角 \(\beta\) 的终边关于 \(y\) 轴对称，所以它们的正弦值相等. 故 \( \sin\beta=\sin\alpha=\frac13 \).
\end{solution}

\begin{problem}
若双曲线 \(x^2-\dfrac{y^2}{m}=1\) 的离心率为 \(\sqrt3\)，则实数 \(m=\fillinblank{}\).

\end{problem}

\begin{answer}
\(2\)
\end{answer}

\begin{solution}
由双曲线方程，\(a^2=1\)，\(b^2=m\)，于是 \(e^2=1+\frac{b^2}{a^2}=1+m\). 又因为 \(e=\sqrt3\)，所以 \(1+m=3\)，解得 \(m=2\).
\end{solution}

\begin{problem}
已知 \(x\ge0\)，\(y\ge0\)，且 \(x+y=1\)，则 \(x^2+y^2\) 的取值范围是\fillinblank{}.

\end{problem}

\begin{answer}
\(\left[\dfrac12,1\right]\)
\end{answer}

\begin{solution}
由 \(x+y=1\)，\(\ x\)，\(y\ge0\)，得 \(x^2+y^2=(x+y)^2-2xy=1-2xy\). 当 \(x=y=\dfrac12\) 时，\(xy\) 取最大值 \(\dfrac14\)，此时 \(x^2+y^2=\dfrac12\)；当 \(x=0\) 或 \(y=0\) 时，\(x^2+y^2=1\). 所以 \(x^2+y^2\) 的取值范围是 \(\left[\dfrac12,1\right]\).
\end{solution}

\begin{problem}
已知点 \(P\) 在圆 \(x^2+y^2=1\) 上，点 \(A\) 的坐标为 \((-2,0)\)，\(O\) 为原点，则 \(\overrightarrow{AO}\cdot\overrightarrow{AP}\) 的最大值为\fillinblank{}.

\end{problem}

\begin{answer}
\(6\)
\end{answer}

\begin{solution}
设 \(P(x,y)\)，则 \(x^2+y^2=1\). \(\overrightarrow{AO}=(2,0)\)，\(\overrightarrow{AP}=(x+2,y)\). 所以 \(\overrightarrow{AO}\cdot\overrightarrow{AP}=2(x+2)=2x+4\). 当 \(x=1\) 时取最大值 \(6\). 故答案为 \(6\).
\end{solution}

\begin{problem}
能够说明“设 \(a\)，\(b\)，\(c\) 是任意实数. 若 \(a>b>c\)，则 \(a+b>c\)”是假命题的一组整数 \(a\)，\(b\)，\(c\) 的值依次为\fillinblank{}.

\end{problem}

\begin{answer}
\(-1\)，\(-2\)，\(-3\)（答案不唯一）
\end{answer}

\begin{solution}
取 \(a=-1\)，\(b=-2\)，\(c=-3\)，则 \(a>b>c\) 成立，但 \(a+b=-3\not>-3\). 所以这一组整数可以说明原命题是假命题.
\end{solution}

\begin{problem}
某学习小组由学生和教师组成，人员构成同时满足以下三个条件：

\begin{circlelist}
\item 男学生人数多于女学生人数；
\item 女学生人数多于教师人数；
\item 教师人数的两倍多于男学生人数.
\end{circlelist}
\begin{enumerate}
\item 若教师人数为 \(4\)，则女学生人数的最大值为\fillinblank{}.
\item 该小组人数的最小值为\fillinblank{}.
\end{enumerate}
\end{problem}

\begin{answer}
（1）\(6\)；（2）\(12\)
\end{answer}

\begin{solution}
设男学生人数为 \(m\)，女学生人数为 \(f\)，教师人数为 \(t\). （1）当 \(t=4\) 时，由条件得 \(m>f>4\)，\(m<8\). 因为 \(m\)，\(f\) 都是整数，所以 \(m\le7\)，从而 \(f\le6\). 故女学生人数最大为 \(6\).

（2）要使总人数最小，只需使 \(t\) 尽可能小. 由 \(m<2t\) 且 \(m>f>t\) 可得，\(t=1\)，\(2\) 时无整数解；当 \(t=3\) 时，可取 \(m=5\)，\(f=4\). 此时总人数最小，为 \(5+4+3=12\). 故答案为 \(6\)；\(12\).
\end{solution}
\section{解答题}

\begin{problem}

已知等差数列 \(\{a_n\}\) 和等比数列 \(\{b_n\}\) 满足 \(a_1=b_1=1\)，\(a_2+a_4=10\)，\(b_2b_4=a_5\).

\begin{enumerate}
\item 求 \(\{a_n\}\) 的通项公式；
\item 求和： \(b_1+b_3+b_5+\cdots+b_{2n-1}\).
\end{enumerate}

\end{problem}

\begin{solution}
（1）设等差数列 \(\{a_n\}\) 的公差为 \(d\). 因为 \(a_2+a_4=10\)，所以 \((a_1+d)+(a_1+3d)=10\). 又 \(a_1=1\)，所以 \(2+4d=10\)，解得 \(d=2\). 所以 \(a_n=a_1+(n-1)d=1+2(n-1)=2n-1\).

（2）设等比数列 \(\{b_n\}\) 的公比为 \(q\). 因为 \(b_2b_4=a_5\)，所以 \((b_1q)(b_1q^3)=a_5\). 由（1）得 \(a_5=9\)，且 \(b_1=1\)，所以 \(q^4=9\). 因为等比数列相邻两项的比值为实数，所以 \(q^2=3\). 于是 \(b_{2k-1}=b_1q^{2k-2}=3^{k-1}\). 所以 \(b_1+b_3+b_5+\cdots+b_{2n-1}=1+3+3^2+\cdots+3^{n-1}=\frac{3^n-1}{2}\).
\end{solution}

\begin{problem}

已知函数 \(f(x)=\sqrt3\cos\left(2x-\dfrac\pi3\right)-2\sin x\cos x\).

\begin{enumerate}
\item 求 \(f(x)\) 的最小正周期；
\item 求证：当 \(x\in\left[-\dfrac\pi4,\dfrac\pi4\right]\) 时，\(f(x)\ge-\dfrac12\).
\end{enumerate}

\end{problem}

\begin{solution}
（1）由 \(f(x)=\sqrt3\cos\left(2x-\frac\pi3\right)-\sin2x\)，得 \(f(x)=\frac{\sqrt3}{2}\cos2x+\frac32\sin2x-\sin2x=\frac{\sqrt3}{2}\cos2x+\frac12\sin2x\). 于是 \(f(x)=\sin\left(2x+\frac\pi3\right)\). 所以 \(f(x)\) 的最小正周期为 \(T=\frac{2\pi}{2}=\pi\).

（2）因为 \(x\in\left[-\frac\pi4,\frac\pi4\right]\)，所以 \(2x+\frac\pi3\in\left[-\frac\pi6,\frac{5\pi}{6}\right]\). 而 \(\sin t\) 在区间 \(\left[-\dfrac\pi6,\dfrac{5\pi}{6}\right]\) 上的最小值为 \(-\dfrac12\)，故 \(f(x)=\sin\left(2x+\frac\pi3\right)\ge-\frac12\).
\end{solution}

\begin{problem}

某大学艺术专业 \(400\) 名学生参加某次测评，根据男女学生人数比例，使用分层抽样的方法从中随机抽取了 \(100\) 名学生，记录他们的分数，将数据分成 \(7\) 组：\([20,30)\)，\([30,40)\)，\(\cdots\)，\([80,90]\)，并整理得到如下频率分布直方图：

\begin{enumerate}
\item 从总体的 \(400\) 名学生中随机抽取一人，估计其分数小于 \(70\) 的概率；
\item 已知样本中分数小于 \(40\) 的学生有 \(5\) 人，试估计总体中分数在区间 \([40,50)\) 内的人数；
\item 已知样本中有一半男生的分数不小于 \(70\)，且样本中分数不小于 \(70\) 的男女生人数相等. 试估计总体中男生和女生人数的比例.
\end{enumerate}

\begin{center}
\texfigure[max width=0.62\linewidth,max totalheight=4.8cm]{img_repaint/2017/beijing_liberal/q17_fig1.tex}
\end{center}
\end{problem}

\begin{solution}
（1）根据频率分布直方图可知，样本中分数不小于 \(70\) 的频率为 \((0.02+0.04)\times10=0.6\)，所以样本中分数小于 \(70\) 的频率为 \(1-0.6=0.4\). 所以从总体的 \(400\) 名学生中随机抽取一人，其分数小于 \(70\) 的概率估计为 \(0.4\).

（2）根据题意，样本中分数不小于 \(50\) 的频率为 \((0.01+0.02+0.04+0.02)\times10=0.9\)，所以样本中分数在区间 \([40,50)\) 内的人数为 \(100-100\times0.9-5=5\). 于是总体中分数在区间 \([40,50)\) 内的人数估计为 \(400\times\frac5{100}=20\).

（3）样本中分数不小于 \(70\) 的学生人数为 \((0.02+0.04)\times10\times100=60\). 又因为样本中分数不小于 \(70\) 的男女生人数相等，所以样本中分数不小于 \(70\) 的男生人数为 \(30\). 于是样本中的男生人数为 \(30\times2=60\)，女生人数为 \(100-60=40\). 所以样本中男生和女生人数的比例为 \(60:40=3:2\)，由分层抽样原理，总体中男生和女生人数的比例估计为 \(3:2\).
\end{solution}

\begin{problem}

如图，在三棱锥 \(P-ABC\) 中，\(PA\perp AB\)，\(PA\perp BC\)，\(AB\perp BC\)，\(PA=AB=BC=2\)，\(D\) 为线段 \(AC\) 的中点，\(E\) 为线段 \(PC\) 上一点.

\begin{enumerate}
\item 求证：\(PA\perp BD\)；
\item 求证：平面 \(BDE\perp\) 平面 \(PAC\)；
\item 当 \(PA //\) 平面 \(BDE\) 时，求三棱锥 \(E-BCD\) 的体积.
\end{enumerate}

\begin{center}
\bitmapinclude[max width=0.58\linewidth,max totalheight=4.8cm]{img/2017/beijing_liberal/q18_fig1.png}
\end{center}
\end{problem}

\begin{solution}
（1）因为 \(PA\perp AB\)，\(PA\perp BC\)，所以 \(PA\perp\) 平面 \(ABC\). 又因为 \(BD\subset\) 平面 \(ABC\)，所以 \(PA\perp BD\).

（2）因为 \(AB=BC\)，且 \(D\) 为 \(AC\) 的中点，所以 \(BD\perp AC\). 由（1）知 \(PA\perp BD\)，而 \(PA\) 与 \(AC\) 都在平面 \(PAC\) 内，且 \(PA\cap AC=A\)，所以 \(BD\perp\) 平面 \(PAC\). 因此平面 \(BDE\perp\) 平面 \(PAC\).

（3）因为 \(PA //\) 平面 \(BDE\)，且平面 \(PAC\cap\) 平面 \(BDE=DE\)，所以 \(PA // DE\). 因为 \(D\) 为 \(AC\) 的中点，所以 \(DE=\frac12PA=1\). 在直角三角形 \(ABC\) 中，\(AB=BC=2\)，所以 \(AC=2\sqrt2\). 又 \(D\) 为 \(AC\) 的中点，故 \(BD=CD=\frac{AC}{2}=\sqrt2\). 由（1）知，\(PA\perp\) 平面 \(ABC\)，所以 \(DE\perp\) 平面 \(PAC\). 且 \(BD\perp CD\)，所以三棱锥 \(E-BCD\) 的体积 \(V=\frac16BD\cdot DC\cdot DE=\frac13\).
\end{solution}

\begin{problem}

已知椭圆 \(C\) 的两个顶点分别为 \(A(-2,0)\)，\(B(2,0)\)，焦点在 \(x\) 轴上，离心率为 \(\dfrac{\sqrt3}{2}\).

\begin{enumerate}
\item 求椭圆 \(C\) 的方程；
\item 点 \(D\) 为 \(x\) 轴上一点，过 \(D\) 作 \(x\) 轴的垂线交椭圆 \(C\) 于不同的两点 \(M\)，\(N\)，过 \(D\) 作 \(AM\) 的垂线交 \(BN\) 于点 \(E\). 求证：\(\triangle BDE\) 与 \(\triangle BDN\) 的面积之比为 \(4:5\).
\end{enumerate}

\end{problem}

\begin{solution}
（1）设椭圆 \(C\) 的方程为 \(\frac{x^2}{a^2}+\frac{y^2}{b^2}=1\)（\(a>b>0\)）. 由题意，\(a=2\)，且 \(e=\frac ca=\frac{\sqrt3}{2}\)，所以 \(c=\sqrt3\). 于是 \(b^2=a^2-c^2=4-3=1\). 故椭圆 \(C\) 的方程为 \(\frac{x^2}{4}+y^2=1\).

（2）设 \(M(m,n)\)，则 \(D(m,0)\)，\(N(m,-n)\)，且 \(m\ne\pm2\)，\(n\ne0\). 直线 \(AM\) 的斜率 \(k_{AM}=\dfrac{n}{m+2}\)，故直线 \(DE\) 的方程为 \(y=-\frac{m+2}{n}(x-m)\). 直线 \(BN\) 的方程为 \(y=\frac{n}{2-m}(x-2)\). 联立
\examdisplaycases{}{
y=-\dfrac{m+2}{n}(x-m),\\[4pt]
y=\dfrac{n}{2-m}(x-2)
}
解得点 \(E\) 的纵坐标 \(y_E=-\frac{n(4-m^2)}{4-m^2+n^2}\). 由点 \(M\) 在椭圆 \(C\) 上，得 \(\frac{m^2}{4}+n^2=1\)，即 \(4-m^2=4n^2\). 所以 \(y_E=-\frac45n\). 又 \(S_{\triangle BDE}=\frac12|BD|\cdot|y_E|=\frac25|BD|\cdot|n|\)，\(S_{\triangle BDN}=\frac12|BD|\cdot|n|\). 所以 \(\triangle BDE\) 与 \(\triangle BDN\) 的面积之比为 \(4:5\).
\end{solution}

\begin{problem}

已知函数 \(f(x)=\e^x\cos x-x\).

\begin{enumerate}
\item 求曲线 \(y=f(x)\) 在点 \((0,f(0))\) 处的切线方程；
\item 求函数 \(f(x)\) 在区间 \(\left[0,\dfrac\pi2\right]\) 上的最大值和最小值.
\end{enumerate}

\end{problem}

\begin{solution}
（1）因为 \(f(x)=\e^x\cos x-x\)，所以 \(f'(x)=\e^x(\cos x-\sin x)-1\). 又 \(f'(0)=0\)，\(\quad f(0)=1\). 所以曲线 \(y=f(x)\) 在点 \((0,f(0))\) 处的切线方程为 \(y=1\).

（2）设 \(h(x)=\e^x(\cos x-\sin x)-1\)，则 \(h'(x)=\e^x(\cos x-\sin x)-\e^x(\sin x+\cos x)=-2\e^x\sin x\). 当 \(x\in\left(0,\dfrac\pi2\right)\) 时，\(h'(x)<0\)，所以 \(h(x)\) 在 \(\left[0,\dfrac\pi2\right]\) 上单调递减. 又 \(h(0)=0\)，因此对任意 \(x\in\left(0,\dfrac\pi2\right]\)，有 \(h(x)<0\)，即 \(f'(x)<0\). 所以 \(f(x)\) 在区间 \(\left[0,\dfrac\pi2\right]\) 上单调递减. 因此 \(f(x)\) 在区间 \(\left[0,\dfrac\pi2\right]\) 上的最大值为 \(f(0)=1\)，最小值为 \(f\left(\dfrac\pi2\right)=-\dfrac\pi2\).
\end{solution}
